
\section{实验结果与分析}
\label{sec:ch3_experiment}

为验证所提 BC-CTDE-MAPPO 方法在天基计算网络协同资源分配场景中的有效性，本节从收敛性、不同负载水平下的系统性能、抗自私行为鲁棒性以及机制有效性四个方面开展仿真实验。实验指标主要包括任务完成率、平均任务时延、Jain 公平指数以及链上信誉演化趋势。为保证对比的完整性，本文选取 MAPPO、MADDPG、Greedy 和 Random 作为基线方法。其中，MAPPO 和 MADDPG 用于比较多智能体强化学习框架的性能差异，Greedy 表示基于局部即时收益的启发式策略，Random 表示随机卸载与资源分配策略。

\subsection{实验设置}
实验场景中包含 20 个异构卫星计算节点，系统共考虑 4 类在轨任务，训练总回合数为 500，每回合包含 200 个时隙。Actor 与 Critic 的学习率分别设置为 $3\times10^{-4}$ 和 $1\times10^{-3}$，折扣因子取 $\gamma=0.99$，GAE 参数取 $\lambda=0.95$，PPO 裁剪系数取 0.2。混合奖励函数中效率项、公平项与可信项的权重分别设置为 0.5、0.3 和 0.2。联盟链默认同步时延设为 1.0 s，任务到达率从 0.3 增加至 1.1，自私节点比例从 0\% 增加至 40\%。详细参数如表~\ref{tab:ch3_params} 所示。

\subsection{收敛性分析}
图~\ref{fig:ch3_conv_comb}(a) 给出了不同训练算法在训练过程中的平均回报收敛曲线。可以看出，随着训练回合增加，各算法的平均回报整体呈现逐步上升并最终趋于稳定的趋势，说明所建多智能体博弈环境与策略优化过程能够有效收敛。其中，BC-CTDE-MAPPO 在约 200 回合后已进入稳定收敛区间，且最终平均回报明显高于其余基线方法；相比之下，MAPPO 的收敛速度略慢，MADDPG 与 IPPO 则表现出更明显的波动和更低的稳态回报。这表明，在引入区块链公共信息板和可信激励机制后，智能体在训练阶段能够更充分地利用系统级上下文，从而降低多智能体环境中的非平稳性。

图~\ref{fig:ch3_conv_comb}(b) 进一步比较了不同算法的收敛回合数。可以看到，BC-CTDE-MAPPO 的收敛回合最少，仅为 210 回合；MAPPO、MADDPG 和 IPPO 分别需要 260、330 和 360 回合才达到稳定状态。结合表~\ref{tab:ch3_convergence} 中的统计结果可知，BC-CTDE-MAPPO 在最终平均回报、收敛速度及后 50 回合回报标准差等指标上均优于对比方法。尤其是在任务到达率 $\lambda=0.9$ 的场景下，其任务完成率达到 92.3\%，明显高于 MAPPO 的 88.1\% 和 MADDPG 的 84.2\%。这说明所提方法不仅具备更快的收敛速度，而且在收敛后仍能保持较高的策略稳定性。

\subsection{不同负载水平下的系统性能分析}
为考察算法在不同业务压力下的适应能力，本文进一步比较了各方法在不同任务到达率下的任务完成率、平均任务时延和 Jain 公平指数，结果分别如图~\ref{fig:ch3_load_comb} 和表~\ref{tab:ch3_load_perf} 所示。

从图~\ref{fig:ch3_load_comb}(a) 可以看出，随着任务到达率由 0.3 增加到 1.1，所有算法的任务完成率均有所下降，这是因为系统负载持续增加后，节点的局部队列长度和资源竞争程度显著提升。然而，BC-CTDE-MAPPO 在整个负载区间内始终保持最优性能，其任务完成率由 98.2\% 下降至 89.1\%，下降幅度明显小于其余基线方法。尤其在高负载场景 $\lambda=1.1$ 下，BC-CTDE-MAPPO 相比 Greedy 和 Random 分别高出 18.3 和 30.2 个百分点，说明所提方法具有更强的任务承载能力。

图~\ref{fig:ch3_load_comb}(b) 给出了不同负载水平下的平均任务时延变化趋势。可以看出，随着任务到达率增大，各算法的任务时延均呈上升趋势，但 BC-CTDE-MAPPO 的增长最缓，在 $\lambda=0.9$ 和 $\lambda=1.1$ 时分别仅为 178 ms 和 214 ms，显著低于 MAPPO、MADDPG 及启发式策略。这说明，所提方法能够通过链上公共摘要感知全局负载状态，并在协同调度时主动规避热点节点，从而降低拥塞传播对系统时延性能的影响。

除效率指标外，本文还使用 Jain 公平指数评价不同算法的负载均衡能力。表~\ref{tab:ch3_load_fairness} 表明，BC-CTDE-MAPPO 在各负载水平下均保持最高公平指数。例如，当 $\lambda=0.9$ 时，其 Jain 公平指数为 0.918，而 MAPPO 和 MADDPG 分别为 0.878 和 0.853，Greedy 与 Random 则仅为 0.781 和 0.731。这说明混合奖励中的公平项能够有效抑制“强节点持续抢占高收益任务”的现象，引导系统逐步形成更均衡的任务分担格局。

\subsection{抗自私行为鲁棒性分析}
为检验所提方法在非合作环境中的鲁棒性，本文设置不同自私节点比例，并比较各算法的任务完成率与 Jain 公平指数，结果如图~\ref{fig:ch3_robust_comb} 和表~\ref{tab:ch3_robustness} 所示。

从图~\ref{fig:ch3_robust_comb}(a) 可以看出，随着自私节点比例从 0\% 提高到 40\%，所有算法的任务完成率均呈下降趋势。这是因为自私节点倾向于拒绝外部卸载或虚报自身状态，从而削弱系统整体协作效率。然而，BC-CTDE-MAPPO 的性能退化最小，其任务完成率仅由 93.2\% 下降至 87.2\%；相比之下，MAPPO 从 89.6\% 下降至 78.4\%，Greedy 和 Random 则下降得更加明显。这表明，所提方法通过区块链的可审计共享和信誉惩罚机制，能够有效缓解自私行为对系统性能的破坏。

图~\ref{fig:ch3_robust_comb}(b) 显示了不同自私节点比例下的 Jain 公平指数变化情况。随着自私节点增多，所有方法的公平指数均有所下降，但 BC-CTDE-MAPPO 始终保持最高水平。例如，在 30\% 自私节点条件下，BC-CTDE-MAPPO 的公平指数仍达到 0.886，而 MAPPO、Greedy 和 Random 分别下降至 0.816、0.722 和 0.659。这说明所提方法不仅能够在效率层面维持更高的任务完成率，也能够在恶劣场景下保持更稳定的全局负载均衡。

\subsection{信誉机制效果与消融实验分析}
图~\ref{fig:ch3_mechanism_comb}(a) 展示了信誉机制下合作节点和自私节点的链上信誉演化过程。可以看出，合作节点的信誉值随着时间步增加逐渐上升，并最终稳定在较高水平；自私节点的信誉值则持续下降。该结果表明，链上信誉反馈机制能够有效区分合作行为与机会主义偏离行为，并将其结果持续回写到后续决策激励中，从而形成“状态感知—策略执行—链上审计—信誉更新”的闭环治理机制。

为了进一步验证各组成模块的贡献，本文在 $\lambda=0.9$ 的场景下开展消融实验，结果如表~\ref{tab:ch3_ablation} 和图~\ref{fig:ch3_mechanism_comb}(b) 所示。与完整 BC-CTDE-MAPPO 相比，去除区块链公共信息板后，任务完成率由 92.3\% 降至 89.1\%，平均任务时延由 178 ms 增加到 194 ms，说明可信共享的系统级上下文对于提升协同调度效果具有关键作用。去除公平奖励项后，系统的 Jain 公平指数从 0.918 降至 0.842，表明公平激励对抑制负载失衡尤为重要。去除可信奖励项后，虽然平均任务时延变化不大，但在 30\% 自私节点场景下的任务完成率下降至 84.1\%，反映出信誉约束在抗自私行为中发挥了不可替代的作用。

\subsection{本节小结}
综上所述，实验结果从多个维度验证了 BC-CTDE-MAPPO 方法的有效性：在收敛性方面，该方法能够以更快速度达到更高稳态回报；在系统性能方面，其在不同负载强度下均表现出更高的任务完成率、更低的平均任务时延和更优的负载均衡能力；在鲁棒性方面，面对不断增加的自私节点比例，所提方法的性能退化幅度显著小于对比算法；在机制层面，区块链公共信息板、公平奖励项与可信奖励项均对最终性能提升具有重要贡献。这些结果表明，所提方法能够在去中心化、自利行为和信息不完全共存的复杂环境中，有效维持天基计算网络的长期稳定协同。
